\documentclass[12pt]{article}
\usepackage[utf8]{inputenc}
\usepackage[T1]{fontenc}
\usepackage{hyperref}
\title{Zulu Language Resources for K-12 Curriculum Development}
\author{Compiled by Claude Sonnet 5.5 (claude-sonnet-5-5) for Andr\'as Kornai}
\date{}
\begin{document}
\maketitle

\section{Introduction}
Zulu (isiZulu; ISO 639-3: zul; Glottocode zulu1248) is an official national language of South Africa with about 16 million speakers. It is used as a language of instruction. Commercial CAPS-aligned materials exist. This search was limited and found mainly publisher and retailer pages. It did not find official Department of Basic Education documents.

\section{Primary Educational Resources}
\begin{itemize}
\item Oxford \textit{Asicije Amakhono} Grade 4 (CAPS) isiZulu workbook: \url{https://oxford.co.za/shop/resources-and-home-learning/workbooks/9780195999815-oxford-asicije-amakhono-ibanga-4-caps/}
\item Oxford \textit{Asicije Amakhono} Grade 5 (CAPS) isiZulu workbook: \url{https://oxford.co.za/shop/school/languages/grade-5-languages/isizulu-grade-5-languages/9780199043774-oxford-asicije-amakhono-ibanga-5-caps/}
\item Twinkl isiZulu topic resources, mainly for the Foundation Phase: \url{https://www.twinkl.com/resources/south-africa-resources/isizulu-resources-south-africa-suid-afrika/topics-isizulu-isixhosa-and-zulu-resources-south-africa-suid-afrika}
\end{itemize}

\section{Linguistic and Reference Materials}
\begin{itemize}
\item isiZulu Handbook and Study Guide, Grades 4--12 (Beryl Lutrin), a reference book covering CAPS language structures and rules: \url{https://shop.snapplify.com/product/the-isizulu-handbook-and-study-guide-grades-4-12}
\item Oxford Bilingual School Dictionary: isiZulu and English, Grades 4--9 (ebook), aimed at additional-language learners: \url{https://oxford.co.za/shop/resources-and-home-learning/resource-book/9780199075935-oxford-bilingual-school-dictionary-isizulu-and-english-grade-4-9-epdf/}
\end{itemize}

\section{Religious and Community Resources}
No religious or community resources were examined in this search.

\section{Digital and Audio Resources}
The ebook and digital items above are the only digital resources found. No audio, corpus or software resources were found in this search.

\section{Curriculum Development Context}
Zulu is already a language of instruction, so a curriculum can build on the existing CAPS home-language framework. The resources found are commercial, so licensing and cost need checking. Official CAPS documents and DBE materials still need to be located, since this search did not find them.
\end{document}